\documentclass[10pt]{article}
\usepackage[a4paper,margin=0.76in]{geometry}
\usepackage{amsmath,amssymb,booktabs,microtype}
\usepackage[colorlinks=true,linkcolor=blue,urlcolor=blue]{hyperref}
\usepackage{enumitem}
\setlist{nosep,leftmargin=*}
\title{Position-Aware Mask-Energy Normalization for a Finite\
Masked Twin-Prime Shell Operator}
\author{Liang Wang\\
School of Mathematics and Statistics\\
Huazhong University of Science and Technology (HUST)\\
Wuhan, China}
\date{3 September 2026}

\newcommand{\norm}[1]{\lVert #1\rVert_2}
\newcommand{\R}{\mathbb{R}}
\newcommand{\kappaout}{\kappa}

\begin{document}
\maketitle

\begin{abstract}
The preceding finite source/operator audit transferred positive operator
polarization to higher origins but showed a downward movement of its all-plus
minimum floor.  We test a predeclared, response-independent position-aware
normalization.  For each unsigned prime component $B_p$ we form the diagonal
mask-energy $G_u=\sum_{p,t}B_p(u,t)^2$ and apply the symmetric congruence
$A^\sharp=D_G^{-1/2}AD_G^{-1/2}$.  The construction uses no source vector,
sign law, or observed response.  On three frozen panels and 648 law-level
rows, the normalized all-plus minimum drop from the low-origin parent to the
higher-origin parent is $0.0262369882$, versus $0.0421511462$ without
normalization, a finite reduction fraction of $0.3775498289$.  The repair is
not uniform: the all-plus mean drop becomes larger, and a fresh mod-4 row is
still negatively aligned.  We therefore record a finite partial repair and a
sharp obstruction, not an asymptotic arithmetic estimate or a twin-prime
theorem.
\end{abstract}

\section{Question and scope}

TPC-353 attached the finite V59 residual
$\beta(t)=\Lambda(t+2)-b^{(2)}(t)$ to a literal two-endpoint
divisibility-masked shell operator.  TPC-354 moved the same construction to
disjoint higher origins.  Its positive alignment census survived, while the
all-plus minimum and mean decreased relative to the low-origin parent.

The present question is deliberately narrower: can a normalization based only
on the position-dependent unsigned mask geometry account for that minimum-floor
movement?  The answer is tested on the locked TPC-353 panel, the locked
TPC-354 panel, and a fresh origins-only panel.  Every conclusion below is
classified as either an exact finite identity, a declared finite model
statement, or a numerically certified finite observation.  No statement here
supplies a growing bound in $X$.

\section{Finite operator and normalization}

Let $I$ be a finite integer interval and let
$S_Q=\{p: p\text{ prime},\ Q<p\leq 2Q\}$.  For exponent $s$ and height $H$,
define the unsigned component
\begin{equation}
 B_p(u,t)=1_{u\ne t}1_{p\nmid u}1_{p\nmid t}
 p\frac{H^{2s}}{(H^2+(u-t)^2)^s}
 \left(1_{u\equiv t\pmod p}-\frac1{p-1}\right).
 \label{eq:block}
\end{equation}
For a predeclared sign law $e=(e_p)_{p\in S_Q}$, the literal operator is
\begin{equation}
 A_e=\sum_{p\in S_Q}e_pB_p.
 \label{eq:operator}
\end{equation}
The endpoint masks in \eqref{eq:block} are part of the definition.

The position-aware geometry energy is
\begin{equation}
 G_u=\sum_{p\in S_Q}\sum_{t\in I}B_p(u,t)^2,
 \qquad D_G=\operatorname{diag}(G_u),
 \qquad A_e^\sharp=D_G^{-1/2}A_eD_G^{-1/2}.
 \label{eq:normalization}
\end{equation}
The square is taken before signs are inserted.  Consequently $G_u$ is
independent of $e$, and it also does not use $\Lambda$, $b$, $\beta$, or an
output vector.  On every audited row $G_u>0$, so the congruence is well
defined.  It is a finite preconditioner, not a claimed uniformly bounded
operator family.

\section{Exact finite identities}

For either $T=A_e$ or $T=A_e^\sharp$ and $\beta=L-b$, finite bilinearity gives
\begin{equation}
 \norm{T\beta}^2=\norm{TL}^2+\norm{Tb}^2-2\langle TL,Tb\rangle.
 \label{eq:polar}
\end{equation}
Writing $E_L=\norm{TL}^2$, $E_b=\norm{Tb}^2$ and
\begin{equation}
 \kappa_T=\frac{2\langle TL,Tb\rangle}{E_L+E_b},
 \qquad R_T=\frac{\norm{T\beta}^2}{E_L+E_b},
 \label{eq:kappa}
\end{equation}
we have the exact finite relation $R_T=1-\kappa_T$.  Cauchy--Schwarz gives
\begin{equation}
 \frac{(\sqrt{E_L}-\sqrt{E_b})^2}{E_L+E_b}
 \leq R_T\leq
 \frac{(\sqrt{E_L}+\sqrt{E_b})^2}{E_L+E_b}.
 \label{eq:cauchy}
\end{equation}
The diagonal congruence changes the finite geometry on which these identities
are evaluated; it does not change their logical status.

\section{Frozen protocol and audit}

The three panels are shown in Table~\ref{tab:protocol}.  Counts, shells,
exponents, height, source cutoff, source convention, and four sign laws are
identical across panels.  The fresh panel was fixed before its responses were
computed.

\begin{table}[h]
\centering\small
\caption{Frozen three-panel protocol.}
\label{tab:protocol}
\begin{tabular}{ll}
\toprule
panel & origins\\
\midrule
low parent (TPC-353) & $6001,8001,10001$\\
higher parent (TPC-354) & $21001,23001,25001$\\
fresh holdout & $29001,33001,37001$\\
\midrule
source counts & $256,512,1024$\\
shell anchors & $Q=24,54,80$\\
kernel exponents & $s=1,2$\\
height / source cutoff & $H=66$ / $50000$\\
sign laws & all-plus, alternating-index, mod-4, half-split\\
\bottomrule
\end{tabular}
\end{table}

The Cartesian product has $3\cdot3\cdot3\cdot2\cdot4=216$ rows per panel,
648 rows in total.  The producer uses the finite V59 midpoint convention and
accumulates shells in increasing order.  A separate checker reconstructs the
source and accumulates shells in reverse order; it does not import the
producer.  Both raw and normalized metrics are compared.  An exact rational
fourteen-point anchor on $[29001,29014]$ with shell $\{5,7\}$ verifies the
unscaled polarization identity, while positivity of the rational geometry
diagonal is checked independently.  Ten in-memory mutations are required to
be rejected by the certificate stress suite.

\section{Results}

Table~\ref{tab:floor} records the all-plus minima and means.  The normalization
reduces the low-to-higher minimum drop, but does not reduce the corresponding
mean drop.

\begin{table}[h]
\centering\small
\caption{All-plus output coefficient $\kappa_T$ by panel.}
\label{tab:floor}
\begin{tabular}{lrrrrrr}
\toprule
 & \multicolumn{3}{c}{raw $A_e$} & \multicolumn{3}{c}{normalized $A_e^\sharp$}\\
panel & min & mean & max & min & mean & max\\
\midrule
low & .692912 & .895612 & .996268 & .690971 & .899454 & .996303\\
higher & .650760 & .874362 & .991350 & .664734 & .874614 & .989806\\
fresh & .654458 & .880122 & .994391 & .664140 & .886090 & .993322\\
\bottomrule
\end{tabular}
\end{table}

The raw higher-panel minimum drop is
\[
  0.69291151430780062-0.65076036812307647
  =0.042151146184724153.
\]
The normalized drop is
\[
  0.69097110464200440-0.66473411648923819
  =0.026236988152766205,
\]
so the descriptive drop-reduction fraction is
\[
  1-\frac{0.026236988152766205}{0.042151146184724153}
  =0.37754982894688971.
\]
On the fresh panel, the normalized minimum is $0.6641398063$, only
$-0.0005943102$ below the higher-panel normalized minimum.  This is evidence
for a finite partial stabilization of the minimum, not a uniform lower-floor
statement.

The mean gives the opposing signal.  The raw low-to-higher mean drop is
$0.021249745559872912$, whereas the normalized mean drop is
$0.024839744603963321$.  For the four laws, the fresh panel retains one
negative mod-4 row in each metric family: the raw negative row is
$(33001,256,Q=24,s=1)$ and the normalized negative row is
$(33001,512,Q=24,s=1)$.  The higher-panel half-split minimum also decreases
under normalization, from $0.0393482615$ to $0.0353891510$.

\begin{table}[h]
\centering\small
\caption{Alignment census over all 648 rows.}
\label{tab:census}
\begin{tabular}{lrrr}
\toprule
metric & positive & negative & unresolved\\
\midrule
raw $A_e$ & 647 & 1 & 0\\
normalized $A_e^\sharp$ & 647 & 1 & 0\\
\bottomrule
\end{tabular}
\end{table}

Thus the normalization improves a selected all-plus minimum statistic while
leaving a law-level sign obstruction and a mean-transfer obstruction intact.

\section{Interpretation and route status}

The reusable structure is a source-independent geometry layer that can be
inserted between a literal masked operator and the finite polarization
interface.  It makes the location of the finite obstruction more explicit:
normalizing unsigned endpoint/mask energy can alter floor behavior without
selecting a canonical sign law.

The strongest licensed claims are:
\begin{itemize}
\item the geometry diagonal and diagonal congruence are exact finite declared
  constructions;
\item polarization and its Cauchy envelope hold exactly for both finite
  operators;
\item the raw/normalized three-panel replay is numerically certified on 648
  rows, with independent reverse-shell and mutation controls;
\item the all-plus minimum-floor reduction is a finite partial repair, while
  mean repair and law-uniform alignment are scoped refuted.
\end{itemize}

No source-uniform arithmetic $L^2$ estimate, growing geometry bound,
canonical sign law, fixed-power credit, Route-B reassembly, or twin-prime
conclusion follows.  The Session-named official Route-A/Route-B evaluator
files are absent from this checkout; the local Bridge-B result is therefore
fail-closed fallback evidence only.

\section{Reproducibility}

The project directory contains the source, independent checker, stress suite,
canonical JSON certificate, proof and theorem ledgers, and this manuscript.
The top-level README gives the exact normal and optimized commands.  The
certificate status is
\begin{center}
\texttt{NUMERICALLY\_CERTIFIED\_FINITE\_POSITION\_AWARE\_MASK\_ENERGY\_NORMALIZATION\_AUDIT}.
\end{center}
The arithmetic advance is `NO', fixed-power credit is zero, and the next
natural experiment is an adversarial position/origin holdout.

\begin{thebibliography}{9}
\bibitem{hardyLittlewood}
G.~H.~Hardy and J.~E.~Littlewood,
Some problems of ``Partitio Numerorum'' III: On the expression of a number as
 a sum of primes, \emph{Acta Math.} 44 (1923), 1--70.
\bibitem{greaves}
G.~Greaves, \emph{Sieves in Number Theory}, Springer, 2001.
\end{thebibliography}

\end{document}
